%% kapa: exam=1 round=22 year=2011 period=1 form=A subject=economics subject_ko=경제원론 seq=20 number=60
%% source: Q-Net 감정평가사 기출문제 (cst00309) artlSeq=5036790 file=extracted/1008/5036790/2011년 제22회 감정평가사 1차시험 기출문제/1교시.hwp
%% conversion: kordoc 4.24.51 -> tools/gen1.py (auto); review=proofread-sonnet; flags=none
\begin{question}{60}
필립스곡선에 관한 설명으로 옳은 것만을 모두 고른 것은?
\begin{framedtext}
ㄱ. 합리적 기대이론에 따르면 기대 인플레이션율이 0\%인 경우에만 단기 필립스곡선은 수직이 된다. \par
ㄴ. 자연실업률가설에 따르면 통화정책에 의해서 장기적으로 자연실업률을 변화시킬 수 있다. \par
ㄷ. 적응적 기대가설 하에서 정부의 재량적 안정화정책은 단기적으로 실업률을 낮출 수 있다. \par
ㄹ. 자연실업률가설에 따르면 장기 필립스곡선은 수직이다.
\end{framedtext}
\begin{choices}
  \item ㄱ, ㄴ
  \item ㄴ, ㄷ
  \item ㄴ, ㄹ
  \item ㄱ, ㄹ
  \item ㄷ, ㄹ
\end{choices}
\end{question}
\answer{5}
